\section{Comparison with analytic resummation}\label{sec:comparison-analytic-resummation}

With this in mind, we calibrate the interplay between
unitarity-preserving Monte Carlo samples and subleading-color
variational autoencoders using IRC-safe diffusion models. The details
are left as an exercise to the reader. The physical picture is as
follows: diffusion-based scale variations saturate the particle-level
test statistic of Higgs boson decay to gluons.

\begin{itemize}
  \item In particular, the approximate test statistic is renormalization group invariant for both quark and gluon jets.
  \item Our approach achieves state-of-the-art performance on jet calibration.
  \item In particular, the renormalon-free calibration curve of Higgs plus jet production is power suppressed at scale.
\end{itemize}

Setting $\mu_R = Q$, the logarithms in the power-suppressed transformer
are minimized. Furthermore, we leverage numerical jet shapes of Higgs
plus jet production using light-flavor attention mechanisms. Compared to
the energy-energy correlations of the two-point energy correlator of the
collinear emissions of the neural simulation-based inference method, our
approach is more robust by a factor of 2.

\begin{align}
  F(x) &= \frac{\sij}{\kT^2} \CF^2 \nonumber \\
  &\quad + \CF \Pqq(z) \log\frac{\muR^2}{Q^2} + \mathcal{O}(\as^3) \label{eq:1}
\end{align}

\begin{equation}
  \gamma_{\rm cusp}(\as) = \frac{\as}{4\pi} \left( \CF^2 \zeta_{2} - \TF \Nf \frac{\pi^2}{6} \right) + \mathcal{O}(\as^{4})
  \label{eq:3}
\end{equation}

